% \label{gen_inst}
\label{sec:method}
\begin{figure}
  \centering
  % \fbox{\rule[-.5cm]{0cm}{4cm} \rule[-.5cm]{4cm}{0cm}}
  \includegraphics[width=1\linewidth]{Fig/Fig1.png}
  \caption{(A) Overall framework for EEG-based object recognition. 
  During training, image-EEG pairs are processed by an image encoder (pre-trained) and an EEG encoder. 
  The objective is to increase the similarity between matched pairs while decreasing it for unmatched pairs. 
  During testing, a few unseen images of target concepts (classes) are processed in advance into templates. 
  Then, the model obtains results by matching test data to templates.
  (B) Architecture of the EEG encoder. Temporal-spatial convolution is used with spatial modules, made with self and graph attention, to reveal spatial features of brain activity. The linear layer is used to project the feature dimension.}
  \label{fig:1}
\end{figure}
\subsection{Overview}
We propose a self-supervised framework, Natural Image Contrast EEG (NICE), to decode images from EEG signals.
The overall framework is depicted in Fig.~\ref{fig:1}(A). 

During training, stimulus-response pairs comprising images and EEG signals are fed into the framework. 
The image and EEG encoder extract features from their respective modalities. 
Contrastive learning is employed to align these features, optimizing the similarity between matched pairs and reducing it for unmatched pairs.
This self-supervised approach enables the model to learn shared information between the two modalities instead of traditional supervised learning with predefined labels. 
Before the test process, we use a few images that belong to the image concepts (classes) to be classified, and have not appeared in the EEG collection experiments.
These images are processed and averaged to obtain one template for each concept.
Next, the EEG encoder processes the test signals, and their similarity with all templates is calculated for matching results.

For the EEG encoder, we adapt temporal-spatial convolution (TSConv), which is widely used in EEG analysis, as shown in 
Fig.~\ref{fig:1}(B). 
In addition, we integrate two plug-and-play spatial modules with self-attention and graph attention.
These modules are instrumental in preserving the spatial characteristics of EEG channels, reflecting intrinsic patterns of brain activity.
Note that the image encoder has been pre-trained on other image datasets. 
The leading studies in computer vision help us with larger sample space, when we cannot easily collect numbers of brain responses.


% \subsection{Preprocessing}
\subsection{Network Architecture}
\subsubsection{EEG Encoder}
\begin{table}[t]
\begin{center}
\renewcommand\arraystretch{1.3}
\caption{Architecture of EEG encoder with SA or GA followed by TSConv}  
\label{tab:1} 
\resizebox{0.7\linewidth}{!}{
\begin{tabular}{lccccc}    
\toprule    
Layer & In & Out & Kernel & Stride & Dimension \\    
\midrule 
% \cmidrule(r){1-2} \cmidrule(r){3-5} 
% \multirow{4}*{2a}
SA \& GA & \multicolumn{5}{c}{\([(b, 1, C, T) \rightarrow (b, 1, C, T)]\), \(b\) is batch} \\
\cline{1-1}
Temporal Conv & 1 & \(k\) & \((1, m_1)\) & (1, 1) & \((b, k, C, T-m_1+1)\)  \\  
Avg Pooling & \(k\) & \(k\) & \((1, m_2)\) & \((1, s)\) & \((b, k, C, (T-m_1-m_2+1)/s+1)\)\\
Spatial Conv & \(k\) & \(k\) & \((C, 1)\) & \((1, 1)\) & \((b, k, 1, (T-m_1-m_2+1)/s+1)\) \\ 
\cline{1-1}
Flatten\&Linear & \multicolumn{5}{c}{\([k*((T-m_1-m_2+1)/s+1)\)\(\rightarrow\) shape of image feature\(]\)}   \\ 
\bottomrule   
\end{tabular}}
\end{center}
\end{table}

Researchers usually arrange the raw EEG trials into two dimensions \begin{math}C\times T\end{math}, in which \begin{math}C\end{math} denotes electrode channels, and \begin{math}T\end{math} denotes time samples.
Convolution along the two dimensions is widely used in deep learning-based EEG analysis models.
In the same way, we present a very concise network architecture as shown in Table~\ref{tab:1}.
One-dimension convolution is first applied to capture the temporal features with \(k\) kernels of size \((1, m_1)\) and stride of \((1, 1)\).
An average pooling layer with a kernel size of \((1, m_2)\) and stride of \((1, s)\) is introduced to alleviate overfitting and smooth the temporal features.
Another one-dimension convolution is then used for spatial features keeping \(k\) kernels of size \((ch, 1)\) and stride of \((1, 1)\), where \(ch\) usually equals the number of electrodes.
After each convolutional layer, batch normalization and exponential linear units (ELU) activation functions are used for better training and nonlinearity \citep{clevert2016fast}.
Finally, a linear layer is added as a projector to transform the features to the same size as the output of the image encoder.

\subsubsection{Image Encoder}
The image features are extracted with the pairing EEG features for contrastive learning.
There have been some excellent computer vision models that can obtain image features with semantic discrimination.
We prefer directly using models pre-trained on other image datasets to get a large sample space. 
The feature extraction parts of Contrastive Language-Image Pre-training (CLIP) \citep{radford2021learning}, Vision Transfomer (ViT) \citep{dosovitskiy2021image}, and Residual Neural Networks (ResNet) \citep{he2016deepa} are tried separately in our framework for demonstration.
All the images for stimuli can be processed in advance, thus speeding up our model training.

\subsubsection{Contrastive Learning}
The framework is constructed with contrastive learning, shown in Algorithm~\ref{alg:1}.
The outputs of the image and EEG encoder are normalized separately. 
Then, the dot product is used to evaluate the similarity of all image-EEG feature pairs.
A scaled temperature parameter is used to adjust distribution probability.
InfoNCE loss \citep{radford2021learning, oord2019representation} is employed as the objective function to increase the similarities between matched pairs and decrease those between unmatched pairs.
After adequate training, the EEG encoder can extract representations similar to corresponding images.
The self-supervised strategy allows us to learn inherent patterns from EEG signals without labels, rather than directly separate different classes with supervised learning. 

\begin{algorithm}[h]
\caption{Natural Image Contrast EEG framework} 
\label{alg:1}
\small
{\fontfamily{cmr}\selectfont
\begin{algorithmic}[1]
    \State \textbf{Input}: (Image, EEG) - stimulus \& response \
    \State \textbf{Model}: Enc\_img - CLIP \&  Enc\_eeg - TSConv 
    \Statex 
    \State \# E - \((batch, channel, electrode, sample)\) \Comment{batch of input EEG}
    \State \# I - \((batch, channel, height, width)\) \Comment{batch of input images}
    \State \# \(\tau\) - learned temperature parameter
    \Statex
    \State \# extract normalized representations from the raw image and EEG
    \State E\_f = Norm(Linear(Enc\_eeg(E))) 
    \State I\_f = Norm(Enc\_img(I))  \Comment{can be obtained before training}
    \Statex
    \State \# scaled pairwise cosine similarity 
    \State logits = dot(E\_f, I\_f.t) * exp(\(\tau\)) 
    \Statex
    \State \# symmetric loss function
    \State labels = arange(batch)
    \State loss\_e = cross\_entropy\_loss(logits, labels, axis=0)
    \State loss\_i = cross\_entropy\_loss(logits, labels, axis=1)
    \State loss = (loss\_e + loss\_i) / 2
\end{algorithmic}}
\end{algorithm}

\subsection{Plug-and-Play Module}
\subsubsection{Self-Attention}
% The electrode channel connectivity plays a vital role in reflecting spatial information.
% However, the commonly used arrangement of EEG data corrupts the original correlations.
% % Some methods recover the spatial dimension to two-dimension, yet bring a high computational complexity.
We utilize two approaches, self-attention (SA) \cite{vaswani2017attentiona}, and graph attention (GA) \cite{velickovic2018graph}, to supplement the EEG encoder as "spatial filters" to encapsulate electrode correlations and reflect the spatial dynamics of brain activity.
We use SA on the electrode channels to evaluate the spatial correlations of EEG data.
The input \(x_{in} \in \mathbb{R}^{C\times T}\) is linearly transformed into equal-sized with weight metrics \(W_q\), \(W_k\), \(W_v\):
\begin{equation}
    x'_{in} = \text{Softmax}(\frac{W_q x_{in} \cdot (W_k x_{in})^T}{\sqrt{d}}) \cdot W_v x_{in}
\end{equation}
where \(x'_{in} \in \mathbb{R}^{C\times T}\) denotes the output after the SA module, \(d\), the time length of EEG data, is a scaled factor to accommodate the Softmax function. 

\subsubsection{Graph Attention}
We employ the GA module to update each electrode with all the others, using the implementation from \cite{brody2022how}.
We treat each electrode as a node \(n_i \in \mathbb{R}^{1\times T}\), \(i=1,...,ch\), which has edges to all the other electrodes \(\mathcal{N}_i\).
The process to update one electrode is as follows:
\begin{equation}
    n'_i = \alpha_{i, i} Wn_i + \sum_{j \in \mathcal{N}_i} \alpha_{i, j} Wn_j
\end{equation}
where \(n'_i\) denotes each node after processing, \(\alpha_{i, j}\) is attention coefficients indicating importance of node \(j\) features to node \(i\), and \(W\) is the weight of linear transformation.
Calculate \(\alpha_{i, j}\) as:
\begin{equation}
    \alpha_{i,j} = \frac{\exp(a^{T}\text{LeakyReLU}(W[n_i \, \Vert \, n_j]))}{\sum_{k \in \mathcal{N}_i \cup \{ i \}}\exp(a^{T}\text{LeakyReLU}(W[n_i \, \Vert \, n_k]))}
\end{equation}
where \(a \in \mathbb{R}^{2T}\) denotes the weight of a feedforward layer for attention calculation, \(()^{T}\) is the transposition operator, and \(\Vert\) is the concatenation operator.
LeakyReLU with a slope of 0.2 is used for nonlinearity in calculation.
Residual connection is employed by integrating the input and output of the SA and GA modules separately, contributing to stable training \citep{he2016deepa}.
